
Table~\ref{tab:validation} reports the validation outcome against the gates fixed
\emph{ex ante}, with each metric assessed against the formulation to which it applies. That
attribution is essential when reading the table: the baseline prescribes supply temperature
from the heating curve and does not propagate it along pipes, so downstream temperature
metrics are structurally outside its scope and are marked as such rather than as failures.

\paragraph{Baseline validation}
Three quantities test the formulation that produces every cost result in this paper, and all
three pass. Annual delivered energy is reproduced to \textbf{1.23}\,\% against a \(2\,\%\)
gate. This anchors the energy scale priced by the decomposition. Network loss is the model's
computed difference between generation and delivered demand and is therefore
\emph{model-implied}, not measured: an independent value would require metered supply-side
energy that the site does not record. Delivered-energy agreement fixes the demand scale to
which that loss is referenced, rather than the loss itself.
The source return temperature matches to \textbf{2.08}\,°C within a \(2.5\)\,°C calibration
tolerance and,
more informatively, has a bias of only \(-0.31\)\,°C. The source return sensor is the only
temperature instrument in the network that does \emph{not} sit downstream of a mixing valve,
and the model is essentially unbiased against it. Heat-pump performance remains within its
expected band, with a mean coefficient of performance of \textbf{2.99}.

\paragraph{Temperature field: validation limitation.}
The supply-temperature field belongs to the temperature-propagating formulation, and two
obstacles prevent quantitative validation here. The first is formulational: the baseline does
not propagate supply temperature, so its trunk drop is identically zero compared with a
measured \(8.2\)\,°C, and the far-end and corridor gates do not apply to it. The second is
instrumental: for the propagating formulation, the far-end error is \(9.21\)\,°C, with a bias
of the same magnitude. Mean absolute error and bias coincide to four decimal places, so every
residual has the same sign, and the magnitude lies squarely within the \(5\) to \(15\)\,°C
offset associated with a three-way mixing valve. The corridor comparison shows the same
signature, with a bias of \(-7.73\)\,°C. This is \emph{consistent with} a fixed sensor-location
offset rather than model error. However, the primary and secondary temperatures were not
measured simultaneously, so the available instrumentation cannot distinguish sensor-location
bias from model error, and the temperature field is therefore not independently validated.
Neither obstacle can be removed through a better pipe model, and the first is a deliberate
property of the formulation used for the reported cost results.

This is the concrete form of an argument that runs throughout the paper. A model resolved more
finely than its metering cannot have that additional resolution verified, and a model that
does not represent a quantity cannot be tested on it. The conclusions are therefore anchored
to annual delivered energy, which the measurements resolve, and hence to the model-implied
network loss on which both formulations agree, rather than to a point temperature field that
neither the instrumentation nor the baseline can support.

\paragraph{Scope of the temperature evaluation.}
Two views of the temperature error are consistent. On the ex-ante validation set, comprising
the held-out nodes over the 744-hour winter window for which the gates in
Table~\ref{tab:valtargets} were fixed, the mean supply-temperature error is \(1.32\)\,°C and
the worst-node error is \(2.27\)\,°C, both \emph{within} their respective \(1.5\)\,K and
\(2.5\)\,K gates. The \(9.21\)\,°C figure in Table~\ref{tab:validation} instead comes from a
separate, more conservative diagnostic based on the far-end node over the full annual record,
without selecting a favourable window. The far end operates the mixing valve closest to fully
open. Appendix~\ref{app:per_node_mae} confirms that the earlier selection was not favourable:
the all-node mean of \(1.68\)\,°C lies between the validation-node and calibration-node group
means, and two excluded nodes are the network's best-performing.

\paragraph{Flow, and the per-step balance.}
The source flow has a mean absolute percentage error of \textbf{33.8}\,\%, and disaggregation
by load band explains rather than excuses it. The absolute error is nearly constant across
bands at roughly \(12\) to \(13\,\mathrm{m^3\,h^{-1}}\), so the percentage is governed by
the denominator: the same physical error is \(46\,\%\) of a
\(30.9\,\mathrm{m^3\,h^{-1}}\) mean flow below one-quarter of peak load and \(13\,\%\) of a
\(102.6\,\mathrm{m^3\,h^{-1}}\) mean flow between one-half and three-quarters of peak load.
Low-load hours account for \(60\,\%\) of the year but carry little energy, which is why a
large relative error in those hours is compatible with an annual energy match of
\(1.23\,\%\). The top band is the exception in both directions, with the largest absolute
error of \(25.5\,\mathrm{m^3\,h^{-1}}\) and a relative error of \(23\,\%\), but comprises
only 56 hours; it is reported rather than absorbed into an average. The per-step energy
balance, comprising generation plus net storage against delivered demand and network loss,
closes to machine precision, with an annual residual below \(10^{-3}\,\mathrm{MWh}\) and
every hour within tolerance. This confirms that the model's internal conservation holds
exactly. The annual demand-energy agreement of \(1.23\,\%\) is the separate measured
comparison on which the cost conclusions rest.

A percentage error can conceal either of two different failures: a model that mistakes the
level or one that fails to follow the variation. Comparing first differences separates them
because a fixed offset cancels. The model tracks the flow level with a correlation of
\textbf{0.91} and its day-to-day change with \textbf{0.80}; for demand, the corresponding
values are \textbf{0.93} and \textbf{0.89}. At hourly resolution, the correlation of
differences falls to \(0.24\), which reflects the metering rather than the model: the record
is sampled at fifteen-minute intervals and resampled to hourly resolution, so hour-to-hour
differences are dominated by sampling noise. The aggregate flow error is therefore a
level-and-denominator effect at low load, not a failure to reproduce how the network varies.
That variation is the property on which the dispatch results depend, since cost accumulates
over days and seasons rather than individual hours.

\paragraph{Post-upgrade assets and structural checks.}
The heat pump, electrode boiler, and storage were installed after the measurement record, so
their dispatch is verified against necessary conditions rather than measurements: the
coefficient of performance remains within its plausible band, the storage state of charge
remains within its operating limits, and simultaneous charging and discharging are
prohibited. Two structural checks close the section. The nonlinear reference with its
quadratic constraints deactivated reproduces the baseline objective to better than
\(0.01\,\%\). On the economic-cost basis, the ordering
\(\mathrm{cost}(\CP) \leq \mathrm{cost}(\Lone)\) also holds for every one of the 135 synthetic
instances, the expected direction for a loss-free relaxation, though not an inequality
guaranteed in general. A departure from either check on this test set would indicate a
formulation error rather than a finding.